% GENERATED FILE. Edit canonical lessons, then run npm run generate:books.
% canonical-source-hash: fnv1a64:875a50245163ecf3
% canonical-lessons: ES-C09-oso

\chapter{-oso --- a Shelf of Words You Were Never Taught}
\label{ch:oso}
\hlchaptermodality{61}

\begin{chapteropening}
\textbf{By the end of this chapter:} \emph{I can read an -oso word as English -ous, make it agree, and choose ser or estar for it.}

Read misterioso, make comida deliciosa agree, and pick estar for nervioso.
\end{chapteropening}

\section[-oso]{-oso}

\label{lesson:ES-C09-oso}



\begin{quote}
Name your two describing words. (\emph{Grande}, \emph{cansado}.)

Two is a short shelf. You are about to have a long one.
\end{quote}

\begin{grammarlens}[title={Grammar Lens: full of}]
\begin{quote}
\textbf{Spanish \emph{-oso} is English \emph{-ous}.}
\end{quote}

Latin \textbf{-ōsus} meant \emph{full of}. Stick it on a noun and you get an adjective: \emph{fāma} is fame or report, so \emph{famōsus} is \textbf{full of fame}.

English holds the same suffix, worn down through French to \emph{-ous}. Not a resemblance — the \textbf{same ending}, arriving by two roads, exactly like \emph{-ción} and \emph{-tion}.

\noindent
\begin{tabularx}{\linewidth}{@{}>{\raggedright\arraybackslash}X>{\raggedright\arraybackslash}X@{}}
\toprule
\textbf{Spanish} & \textbf{English} \\
\midrule
\emph{famoso} & \textbf{famous} \\
\emph{delicioso} & \textbf{delicious} \\
\emph{curioso} & \textbf{curious} \\
\emph{precioso} & \textbf{precious} \\
\emph{nervioso} & \textbf{nervous} \\
\bottomrule
\end{tabularx}

You were taught none of those. You can read all of them.
\end{grammarlens}

\begin{grammarlens}[title={Grammar Lens: and you already know what to do with it}]
An \emph{-oso} word is an ordinary describing word, so every rule you just learned applies without being restated:

\begin{itemize}
\raggedright
  \item it goes \textbf{after} the noun — \emph{un profesor famoso}
  \item it ends in \emph{-o}, so it \textbf{shifts} — \emph{famoso} / \emph{famosa}
  \item and the verb depends on what you mean:
\end{itemize}

\noindent
\begin{tabularx}{\linewidth}{@{}>{\raggedright\arraybackslash}X>{\raggedright\arraybackslash}X@{}}
\toprule
\textbf{} & \textbf{} \\
\midrule
\emph{Es famoso.} & He is famous — that is what he is. \\
\emph{Estoy nervioso.} & I am nervous — right now. \\
\bottomrule
\end{tabularx}

That last pair is the whole \emph{ser}/\emph{estar} contrast again, and you did not have to be told which verb to use. You worked it out from the meaning.
\end{grammarlens}

\begin{culture}
\textbf{The ending is reliable. The root still has to be shared.}

\emph{Hermoso} is \emph{beautiful}, not \emph{hermous} — because the root is \emph{formōsus}, from \emph{forma}, and English simply never borrowed that one.

So \emph{-oso} tells you the word is an adjective and roughly how it is built. It only hands you the meaning when English took the same root. That is the same deal every rule in this book has offered: a \textbf{lead}, not a guarantee.
\end{culture}

\subsection*{Your turn}
Say these aloud:

\begin{itemize}
\raggedright
  \item \emph{famoso, delicioso, curioso}
  \item \emph{Estoy nervioso}
\end{itemize}

\emph{Twice through:}

\begin{tcolorbox}[breakable,colback=teal!4,colframe=teal!35!black,title={Before you move on}]
\emph{-oso} in English? (\emph{\textbf{-ous}}.) Does it shift for a woman? (\textbf{Yes} — it ends in \emph{-o}.)

And what does the ending \emph{not} promise? (That English \textbf{took the same root}.)
\end{tcolorbox}
